\chapter{Massless Scalar Electrodynamics and Dimensional Transmutation}
\label{ch:massless-scalar-ed}

In renormalized massless scalar electrodynamics, the scalar field acquires
an arbitrary vacuum expectation value $\langle\phi\rangle$ even though
classically it vanishes. This means $\langle\phi\rangle$ emerges as a
spontaneously generated physical mass scale, with both scalar and vector
particles developing dynamical masses proportional to it. This phenomenon
is called \emph{dimensional transmutation}: a theory with no classical
mass scale develops a preferred mass scale through radiative corrections.

\section{The Lagrangian and Unitary Gauge}

The Lagrangian of massless scalar electrodynamics is
\begin{equation}
    \mathcal{L} = -\frac{1}{4}F_{\mu\nu}F^{\mu\nu}
    + \left[(\partial_\mu - ieA_\mu)\phi^*\right]
    \left[(\partial^\mu + ieA^\mu)\phi\right]
    - \frac{\lambda_0}{6}(\phi^*\phi)^2,
\end{equation}
where $e$ is the renormalized coupling constant and $\lambda_0$ is the
bare scalar coupling, necessary to renormalize the scalar-scalar
scattering amplitude. Transforming away the phase of $\phi$ by going to
unitary gauge and rescaling $\phi\to\phi/\sqrt{2}$, the Lagrangian
becomes
\begin{equation}
    \mathcal{L} = -\frac{1}{4}F_{\mu\nu}F^{\mu\nu}
    + \frac{1}{2}e^2A_\mu A^\mu\phi^2
    + \frac{1}{2}\partial_\mu\phi\,\partial^\mu\phi
    - \frac{\lambda_0}{4!}\phi^4,
\end{equation}
subject to the Lorenz condition $\partial_\mu A^\mu = 0$. The
corresponding action is
\begin{equation}
    S[A,\phi] = \int d^4x\left[\frac{1}{2}A^\mu
    (\partial^2 + e^2\phi^2)A_\mu
    - \frac{1}{2}\phi\,\partial^2\phi
    - \frac{\lambda_0}{4!}\phi^4\right].
\end{equation}

\section{Integrating Out the Vector Field}

To obtain an effective action for $\phi$ alone, we integrate out the
vector field,
\begin{equation}
    \exp iS[\phi] = \mathcal{N}\int(\mathcal{D}A)\,
    \delta[\partial_\mu A^\mu]\exp iS[A,\phi],
\end{equation}
where $\delta[\partial_\mu A^\mu]$ enforces the Lorenz constraint. Each
transverse component of $A_\mu$ contributes a Gaussian determinant,
\begin{align}
    \left[\det(\partial^2 + e^2\phi^2)\right]^{-\frac{1}{2}}
    &= \exp\left[-\frac{1}{2}\ln\det(\partial^2 + e^2\phi^2)\right]
    \notag\\
    &= \exp\left[-\frac{1}{2}\mathrm{Tr}\ln(\partial^2
    + e^2\phi^2)\right].
\end{align}
There are three such transverse components, so the effective action
becomes
\begin{equation}
    S[\phi] = -\int d^4x\left(\frac{1}{2}\phi\,\partial^2\phi
    + \frac{\lambda_0}{4!}\phi^4\right)
    - \frac{3}{2}\mathrm{Tr}\ln(\partial^2 + e^2\phi^2).
\end{equation}
The one-loop effective potential for $\phi$, to lowest order in
$\lambda_0$, is then
\begin{equation}
    U(\rho) = \frac{\lambda_0}{4!}\rho^4
    + \frac{3}{2}\int\frac{d^4k_E}{(2\pi)^4}
    \ln(k_E^2 + e^2\rho^2),
\end{equation}
where the second term sums all one-loop graphs with external $\phi$ lines
and vector boson running in the loop, since we integrated out $A_\mu$.
Scalar loops contribute only at higher order.

\section{The Renormalized Effective Potential}

Using the machinery developed in the previous chapter, specifically
Eq.\eqref{eq:effective_potential}, the renormalized effective potential
in this case takes the form
\begin{equation}
    U(\rho) = \rho^4\left[\frac{\lambda}{4!}
    + \frac{3e^4}{64\pi^2}\left(-\frac{1}{2}
    + \ln\frac{e^2\rho^2}{\mu^2}\right)\right],
    \label{eq:renormalized_eff_potential}
\end{equation}
where $\lambda$ is the running coupling constant at the renormalization
point $\mu$. Note that $e^4$ is a one-loop result and $\lambda$ is a
tree-level result, but they can be comparable since they are independent
coupling constants. This is what allows the cancellation that was
impossible in the pure scalar case.

\section{Spontaneous Symmetry Breaking and Particle Masses}

Differentiating \eqref{eq:renormalized_eff_potential},
\begin{align}
    U'(\rho) &= 4\rho^3\left[\cdots\right]
    + \rho^4\cdot\frac{3e^4}{64\pi^2}\cdot
    \frac{1}{\frac{e^2\rho^2}{\mu^2}}\cdot\frac{2\rho e^2}{\mu^2}
    \notag\\
    &= 4\rho^3\left[\cdots\right]
    + \frac{3e^4}{64\pi^2}\rho^3\cdot 2 \notag\\
    &= \rho^3\left[4\cdot\frac{\lambda}{4!}
    + 4\cdot\frac{3e^4}{64\pi^2}\left(-\frac{1}{2}\right)
    + 4\cdot\frac{3e^4}{64\pi^2}\ln\frac{e^2\rho^2}{\mu^2}
    + \frac{3e^4}{64\pi^2}\cdot 2\right],
\end{align}
\begin{equation}
    \Rightarrow U'(\rho) = 4\rho^3\left[\frac{\lambda}{4!}
    + \frac{3e^4}{64\pi^2}\ln\frac{e^2\rho^2}{\mu^2}\right].
\end{equation}
Setting $U'(\langle\phi\rangle) = 0$ for a nontrivial root,
\begin{align}
    \frac{\lambda}{4!}
    + \frac{3e^4}{64\pi^2}\ln\frac{e^2\langle\phi\rangle^2}{\mu^2}
    &= 0 \notag\\
    \Rightarrow\quad
    \frac{3e^4}{64\pi^2}\ln\frac{e^2\langle\phi\rangle^2}{\mu^2}
    &= -\frac{\lambda}{4!}.
\end{align}
Unlike the pure scalar case, this root is within the domain of validity
of the approximation since $\lambda$ and $e^4$ are independent and can
be chosen to make the two terms comparable without requiring the
logarithm to be large. Using this condition to eliminate $\mu$ from
\eqref{eq:renormalized_eff_potential},
\begin{align}
    U(\rho) &= \rho^4\left[-\frac{3e^4}{64\pi^2}
    \ln\frac{e^2\langle\phi\rangle^2}{\mu^2}
    + \frac{3e^4}{64\pi^2}\left(-\frac{1}{2}\right)
    + \frac{3e^4}{64\pi^2}\ln\frac{e^2\rho^2}{\mu^2}\right]
    \notag\\
    &= \frac{3e^4\rho^4}{64\pi^2}\left[-\frac{1}{2}
    + \ln\frac{\rho^2}{\langle\phi\rangle^2}\right].
\end{align}
The scalar particle mass follows from $m_S^2 \equiv U''(\langle\phi\rangle)$.
Setting $a\equiv 3e^4/(64\pi^2)$ for brevity, we differentiate
$U(\rho) = a\rho^4\left[-\frac{1}{2}+\ln\frac{\rho^2}{\langle\phi\rangle^2}\right]$,
\begin{align}
    U'(\rho) &= a\cdot 4\rho^3\left[-\frac{1}{2}
    + \ln\frac{\rho^2}{\langle\phi\rangle^2}\right]
    + a\rho^4\cdot\frac{\langle\phi\rangle^2}{\rho^2}
    \cdot\frac{2\rho}{\langle\phi\rangle^2} \notag\\
    &= 2a\rho^3\left[2\left(-\frac{1}{2}
    + \ln\frac{\rho^2}{\langle\phi\rangle^2}\right) + 1\right]
    \notag\\
    &= 4a\rho^3\ln\frac{\rho^2}{\langle\phi\rangle^2}.
\end{align}
Differentiating once more,
\begin{align}
    U''(\rho) &= 12a\rho^2\ln\frac{\rho^2}{\langle\phi\rangle^2}
    + 4a\rho^3\cdot\frac{\langle\phi\rangle^2}{\rho^2}
    \cdot\frac{2\rho}{\langle\phi\rangle^2} \notag\\
    &= 12a\rho^2\ln\frac{\rho^2}{\langle\phi\rangle^2} + 8a\rho^2.
\end{align}
Evaluating at $\rho = \langle\phi\rangle$,
\begin{align}
    U''(\langle\phi\rangle)
    &= 12a\langle\phi\rangle^2\ln\frac{\langle\phi\rangle^2}
    {\langle\phi\rangle^2} + 8a\langle\phi\rangle^2 \notag\\
    &= 8\cdot\frac{3e^4}{64\pi^2}\langle\phi\rangle^2,
\end{align}
\begin{equation}
    \Rightarrow\quad\boxed{m_S^2 = \frac{3e^4}{8\pi^2}
    \langle\phi\rangle^2.}
\end{equation}
The vector field acquires mass through the Higgs mechanism,
\begin{equation}
    m_V^2 = e^2\langle\phi\rangle^2.
\end{equation}
Both masses are proportional to $\langle\phi\rangle$, confirming that
the dynamically generated VEV sets the mass scale of the theory.

\section{Dimensional Transmutation}

To renormalize the theory we introduce a momentum cutoff, which
represents a hidden scale. After renormalization this hidden scale is
replaced by an arbitrary finite renormalization point $\mu$. In a
classically massless theory one might expect all choices of $\mu$ to be
equivalent, but what we have found is that there is a preferred
renormalization point: the one at which the particles develop mass
dynamically. Dimensional transmutation refers precisely to this
emergence of a preferred scale in a theory with no intrinsic mass,
whereby the infinite cutoff momentum becomes transmuted into a finite
physical mass. This mechanism is widely believed to operate in QCD with
massless quarks, where no intrinsic mass scale exists in the Lagrangian
yet hadrons acquire mass dynamically.
